\chapter{Requirements and Analysis}

\section{Problem Definition}

\subsection{Challenge}

Hearing-impaired individuals face significant communication barriers in daily life due to:
\begin{enumerate}
    \item Limited number of sign language interpreters
    \item General population lacks sign language proficiency
    \item Communication gaps in education, healthcare, and public services
    \item Lack of accessible assistive technology solutions
\end{enumerate}

\subsection{Solution Approach}

Develop an automated system that:
\begin{itemize}
    \item Recognizes Indian Sign Language gestures in real-time
    \item Converts gestures to text and speech output
    \item Operates without requiring human interpreters
    \item Is accessible and easy to use
    \item Can be deployed on standard computers with webcams
\end{itemize}

\section{Requirements Specification}

\subsection{Functional Requirements}

\begin{enumerate}
    \item \textbf{Hand Detection}: Real-time detection of one or more hands in video feed
    \item \textbf{Gesture Classification}: Accurate classification of recognized gestures
    \item \textbf{Text Conversion}: Convert recognized gestures to text output
    \item \textbf{Speech Synthesis}: Convert text to speech using text-to-speech engine
    \item \textbf{Training Interface}: Dashboard for collecting and labeling training data
    \item \textbf{Testing Interface}: Real-time gesture recognition testing
    \item \textbf{Typing Mode}: Use gestures to type characters with text-to-speech
    \item \textbf{Data Management}: Store and manage training datasets
    \item \textbf{Model Management}: Save and load trained models
\end{enumerate}

\subsection{Non-Functional Requirements}

\begin{enumerate}
    \item \textbf{Performance}: Real-time processing at 30+ FPS
    \item \textbf{Accuracy}: Achieve >85\% accuracy for trained gestures
    \item \textbf{Reliability}: Consistent performance across lighting conditions
    \item \textbf{Usability}: Intuitive interface for non-technical users
    \item \textbf{Portability}: Cross-platform support (Windows, Linux, macOS)
    \item \textbf{Scalability}: Support for adding new gestures without retraining all models
    \item \textbf{Robustness}: Handle multiple hand distances and angles
\end{enumerate}

\section{Planning and Scheduling}

\subsection{Project Phases}

\begin{table}[H]
\centering
\begin{tabular}{|c|c|c|}
\hline
\textbf{Phase} & \textbf{Duration} & \textbf{Deliverables} \\
\hline
Analysis & Week 1-2 & Requirements, Design Documents \\
\hline
Development & Week 3-6 & Core System Implementation \\
\hline
Dataset Creation & Week 4-5 & Training Data Collection \\
\hline
Model Training & Week 5-6 & Trained Models \\
\hline
Testing & Week 7 & Test Results, Bug Fixes \\
\hline
Integration & Week 7 & Dashboard Integration \\
\hline
Documentation & Week 8 & Final Report, User Guide \\
\hline
\end{tabular}
\caption{Project Timeline}
\end{table}

\subsection{Gantt Chart (Month-wise Project Journey)}

\begin{table}[H]
\centering
\small
\begin{tabular}{|p{4.2cm}|c|c|c|c|c|c|}
\hline
\textbf{Phase} & \textbf{July} & \textbf{August} & \textbf{September} & \textbf{October} & \textbf{November} & \textbf{December} \\
\hline
Requirement Analysis & X &  &  &  &  &  \\
\hline
Setup and Environment Configuration &  & X &  &  &  &  \\
\hline
Backend Development &  &  & X &  &  &  \\
\hline
Frontend Development &  &  &  & X &  &  \\
\hline
Integration and Testing &  &  &  &  & X &  \\
\hline
Deployment and Documentation &  &  &  &  &  & X \\
\hline
\end{tabular}
\caption{Gantt Chart for Project Journey (July to December)}
\end{table}

\section{Software and Hardware Requirements}

\subsection{Software Requirements}

\begin{table}[H]
\centering
\begin{tabular}{|p{3.2cm}|p{4cm}|p{2.4cm}|p{4.2cm}|}
\hline
\textbf{Software Component} & \textbf{Minimum Requirement} & \textbf{Recommended} & \textbf{Purpose} \\
\hline
Operating System & Windows 10 / Ubuntu 18.04 / macOS 10.13 & Windows 11 / Ubuntu 22.04 / macOS 12+ & Application runtime and device management \\
\hline
Python & Python 3.8+ & Python 3.10+ & Core language for model, CV, and dashboard modules \\
\hline
Package Manager & pip / conda & pip (latest) + virtual environment & Dependency installation and environment isolation \\
\hline
IDE/Editor & Any text editor & VS Code / PyCharm & Development, debugging, and code maintenance \\
\hline
Libraries & OpenCV, MediaPipe, TensorFlow, Streamlit, pyttsx3 & Latest stable releases & Vision pipeline, inference, UI, and speech output \\
\hline
Version Control & Git (optional) & Git with remote repository & Source tracking, backup, and collaboration \\
\hline
\end{tabular}
\caption{Software Requirements Specification}
\end{table}

\subsection{Hardware Requirements}

\begin{table}[H]
\centering
\begin{tabular}{|p{3.2cm}|p{3.8cm}|p{3.2cm}|p{3.4cm}|}
\hline
\textbf{Component} & \textbf{Minimum Specification} & \textbf{Recommended} & \textbf{Purpose} \\
\hline
Processor & Intel i5 / AMD Ryzen 5 & Intel i7 / Ryzen 7 or better & Real-time frame processing and model inference \\
\hline
RAM & 4 GB & 8 GB or above & Smooth multi-module execution \\
\hline
Storage & 2 GB free & 10 GB+ SSD free & Dependencies, model files, logs, and dataset growth \\
\hline
Webcam & 640x480 resolution & 720p or 1080p USB webcam & Gesture capture and landmark detection \\
\hline
Display & 1366x768 & Full HD (1920x1080) & Better dashboard usability and visualization \\
\hline
Audio Output & Basic speaker/headphone & Noise-free speaker/headset & Text-to-speech feedback \\
\hline
\end{tabular}
\caption{Hardware Requirements Specification}
\end{table}

\section{Preliminary Product Description}

The Indian Sign Language (ISL) Detection System is a desktop application consisting of:

\begin{enumerate}
    \item \textbf{Core Engine}
    \begin{itemize}
        \item Hand landmark detection using MediaPipe
        \item Gesture classification using CNN models
        \item Point history tracking for complex gestures
    \end{itemize}
    
    \item \textbf{User Interface}
    \begin{itemize}
        \item Streamlit-based web dashboard
        \item Command-line interface (optional)
        \item Real-time video preview
    \end{itemize}
    
    \item \textbf{Data Management}
    \begin{itemize}
        \item Training data collection and storage
        \item Model serialization and loading
        \item Configuration management
    \end{itemize}
    
    \item \textbf{Output Generation}
    \begin{itemize}
        \item Text output of recognized gestures
        \item Speech synthesis using pyttsx3
        \item Gesture-to-character mapping
    \end{itemize}
\end{enumerate}

\section{Feasibility Study}

\subsection{Technical Feasibility}

\begin{itemize}
    \item \textbf{Positive Factors}
    \begin{itemize}
        \item MediaPipe provides robust hand detection with proven accuracy
        \item TensorFlow/Keras enable fast model development and iteration
        \item Streamlit simplifies web UI development without frontend expertise
        \item Python ecosystem is mature, well-documented, and widely supported
        \item TensorFlow Lite enables efficient model deployment on various devices
    \end{itemize}
    
    \item \textbf{Challenges}
    \begin{itemize}
        \item Real-time video processing at 30+ FPS requires careful optimization
        \item Limited availability of public Indian Sign Language datasets
        \item Complex hand gestures with similar hand shapes require careful feature engineering
        \item Handling environmental variability (lighting, backgrounds, hand sizes)
    \end{itemize}
\end{itemize}

\subsection{Economic Feasibility}

\begin{itemize}
    \item All tools, libraries, and frameworks used are completely open-source with zero licensing cost
    \item Minimal hardware investment required—standard consumer laptops are sufficient
    \item Standard USB webcam (< \$30) suitable for implementation
    \item Development uses freely available IDEs, compilers, and tools (VS Code, Python)
    \item No cloud subscriptions or paid services required for core functionality
\end{itemize}

\subsection{Operational Feasibility}

\begin{itemize}
    \item System is designed for end-users with minimal technical knowledge
    \item Simple one-command installation using pip package manager and bash scripts
    \item Web-based Streamlit dashboard provides intuitive, self-explanatory interface
    \item Can be extended to support new gestures through dashboard without code modifications
    \item Minimal training required for users to operate the system
    \item System documentation and inline code comments support long-term maintenance
\end{itemize}

\subsection{Conclusion}

All feasibility aspects are favorable. The project is technically sound with proven component technologies, economically viable with zero-cost open-source stack, and operationally feasible with proper planning and modular design.

\section{Detailed Use Cases and Scenarios}

\subsection{Primary Use Case: Real-Time Gesture Recognition}

\begin{itemize}
    \item \textbf{Actor}: Hearing-impaired user
    \item \textbf{Precondition}: System is running, camera feed is active
    \item \textbf{Main Flow}:
    \begin{enumerate}
        \item User makes a trained gesture in front of the camera
        \item System detects hand landmarks in real-time
        \item Landmarks are preprocessed and normalized
        \item Gesture is classified using the trained model
        \item Recognized gesture text is displayed on screen
        \item Gesture is spoken aloud using text-to-speech
        \item Gesture appears in history log for reference
    \end{enumerate}
    \item \textbf{Postcondition}: Gesture is successfully recognized and communicated
\end{itemize}

\subsection{Secondary Use Case: Training New Gestures}

\begin{itemize}
    \item \textbf{Actor}: System administrator or power user
    \item \textbf{Precondition}: Streamlit dashboard is open
    \item \textbf{Main Flow}:
    \begin{enumerate}
        \item User selects "Train New Gesture" mode from dashboard
        \item User enters name for new gesture
        \item System captures images of user performing the gesture multiple times
        \item User provides images via webcam and dashboard upload
        \item Images are stored in organized directory structure
        \item Model retraining is triggered
        \item New gesture is added to classifier vocabulary
        \item System is restarted to load updated model
    \end{enumerate}
    \item \textbf{Postcondition}: New gesture is integrated and recognized by system
\end{itemize}

\subsection{Tertiary Use Case: Typing Mode}

\begin{itemize}
    \item \textbf{Actor}: Hearing-impaired user wanting to type text
    \item \textbf{Precondition}: System is in typing mode, gesture-to-character mapping is loaded
    \item \textbf{Main Flow}:
    \begin{enumerate}
        \item User makes gesture corresponding to desired character
        \item System recognizes gesture and maps to character
        \item Character is appended to text buffer on screen
        \item Character is pronounced individually via TTS
        \item User continues gesturing to spell out complete word
        \item User makes special gesture to complete word and trigger full word pronunciation
    \end{enumerate}
    \item \textbf{Postcondition}: Complete word or sentence has been typed and spoken
\end{itemize}

\section{System Constraints and Limitations}

\subsection{Environmental Constraints}

\begin{itemize}
    \item \textbf{Lighting Requirements}: System performs best in well-lit environments (> 300 lux). Performance degrades in low-light or shadow conditions.
    \item \textbf{Background Requirements}: Complex or moving backgrounds may reduce landmark detection accuracy. Static, non-distracting backgrounds are preferred.
    \item \textbf{Distance Requirements}: Hand should be 30-100 cm from camera for optimal detection.
    \item \textbf{Frame Rate}: Minimum 15 FPS camera recommended; processing assumes 30 FPS input for smooth temporal tracking.
\end{itemize}

\subsection{Technical Constraints}

\begin{itemize}
    \item \textbf{Gesture Vocabulary Limit}: Current model supports up to 20-30 distinct gestures with practical training overhead.
    \item \textbf{Hand Occlusion}: Partially hidden or overlapped hands reduce landmark accuracy.
    \item \textbf{Multi-Hand Limitation}: Current implementation handles up to 2 hands simultaneously; more than 2 hands not supported.
    \item \textbf{Processing Latency}: ~100ms latency between gesture and text output; not suitable for ultra-real-time applications.
    \item \textbf{Memory Usage}: Requires 300-500 MB RAM for all components; not suitable for very low-resource systems.
\end{itemize}

\subsection{User Constraint}

\begin{itemize}
    \item \textbf{Gesture Consistency}: Users must perform gestures with reasonable consistency to achieve high recognition rates.
    \item \textbf{Hand Position}: Hands must be visible and unobstructed to camera for detection to work.
    \item \textbf{Gesture Vocabulary}: Only pre-trained gestures will be recognized; spontaneous or untrained gestures will result in "Unknown" classification.
\end{itemize}

\section{Assumptions and Dependencies}

\subsection{Assumptions}

\begin{itemize}
    \item System assumes internet availability is \textbf{not} required for core recognition functionality (all models are loaded locally)
    \item Users are assumed to perform gestures deliberately and consciously (not unintentional hand movements)
    \item Training data is assumed to be representative of actual usage contexts
    \item Audio output hardware (speakers) is assumed to be available for text-to-speech
    \item Users are assumed to have basic computer literacy (able to start applications and use dashboard)
\end{itemize}

\subsection{Dependencies}

\begin{itemize}
    \item \textbf{External Libraries}: All dependencies (MediaPipe, TensorFlow, OpenCV, Streamlit, pyttsx3) must be installed and functioning
    \item \textbf{Model Files}: Pre-trained TFLite model files must be present in correct directories
    \item \textbf{Label Files}: CSV files mapping gesture class indices to names must be available
    \item \textbf{Python Version}: Python 3.8+ required for language features and library compatibility
    \item \textbf{Operating System}: Webcam drivers and audio subsystem must be functional
\end{itemize}

\section{Data Flow and Information Requirements}

\subsection{Input Data}

\begin{itemize}
    \item \textbf{Video Stream}: Continuous RGB frames from webcam (640x480, 30 FPS)
    \item \textbf{Landmark Coordinates}: 21 hand landmarks per hand (x, y, z, confidence) from MediaPipe
    \item \textbf{Preprocessed Features}: Normalized 42-element feature vector per hand
    \item \textbf{Training Data}: Labeled gesture images for model training
\end{itemize}

\subsection{Processing Data}

\begin{itemize}
    \item \textbf{Feature Vector}: Normalized landmark coordinates stored as 42-element numpy arrays
    \item \textbf{Model State}: Intermediate activations in neural network layers during inference
    \item \textbf{History Buffer}: Deque of recent gesture predictions for temporal filtering
\end{itemize}

\subsection{Output Data}

\begin{itemize}
    \item \textbf{Gesture Label}: Recognized gesture name as string
    \item \textbf{Confidence Score}: Probability of predicted gesture class
    \item \textbf{Text Output}: Displayed gesture name on screen
    \item \textbf{Audio Output}: Synthesized speech output via speakers
    \item \textbf{Gesture History}: Log of recognized gestures with timestamps
\end{itemize}

\section{Security and Safety Requirements}

\subsection{Data Privacy}

\begin{itemize}
    \item All processing occurs locally on user's machine; no data sent to external servers
    \item Training data stored locally; user has full control over dataset location
    \item No personal information is collected or logged
    \item Video frames are processed and discarded; not stored permanently
\end{itemize}

\subsection{System Safety}

\begin{itemize}
    \item Error handling for missing model files with graceful degradation
    \item Webcam access validation to prevent unauthorized camera usage
    \item Non-blocking speech output to prevent UI freeze
    \item Memory management to prevent buffer overflows or resource exhaustion
\end{itemize}

\section{Acceptance Criteria}

The system is considered acceptable if it meets the following criteria:

\begin{enumerate}
    \item \textbf{Accuracy}: Achieves $\geq 90\%$ recognition accuracy on trained gestures in controlled lighting conditions
    \item \textbf{Performance}: Processes video at $\geq 25$ FPS on standard laptop hardware
    \item \textbf{Usability}: Dashboard interface can be used by non-technical users without training
    \item \textbf{Reliability}: System runs for $\geq 2$ hours continuously without crashes
    \item \textbf{Extensibility}: New gestures can be added and integrated with $< 10$ minutes manual work
    \item \textbf{Documentation}: Complete API documentation and user guide are provided
    \item \textbf{Cross-Platform}: Tested and functional on Windows 10+, Ubuntu 20.04+, and macOS 10.15+
\end{enumerate}
